A query optimizer that rewrites a query wrongly returns a wrong answer without any error, so users have no reason to
doubt it. Random SQL fuzzers rarely reach the narrow conditions under which a single rewrite rule fires, and a bug
found in a released version is often already fixed in the development branch. We present \tool, a framework of four
fuzzing methods for the DuckDB database that run side by side and share one execution judge, one oracle layer and one
triage pipeline: (A) classic generic fuzzing as a baseline, (B) fuzzers aimed at features new in DuckDB~2.0, (C) test
generators written by a large language model (LLM) and guided by new-code hints and fix dossiers, and (D) a
rule-targeted loop. Method D mines a \emph{rule card} for each optimizer rule from its source code and fix history,
asks the LLM for a bug hypothesis and a generator program, and accepts the generator only if executing it shows that
the target rule fired. With these methods we filed 16 issues on DuckDB~2.0.0: 8 silent wrong results, 5 internal
errors and 3 spurious errors. Maintainers have fixed 4 of them, 9 more have open fix pull requests (one of them ours),
and Method D alone found 12. One unattended run of Method D spent 40 LLM calls on 13 targets and produced 9 candidates
that became 6 issues. Generic fuzzing with 70{,}980 SQLsmith queries found no new bug.
